\section{Conflicts between a unit and an independent point}
\label{sec:asymmetric}
\label{sec:asymmetric-test}

We prove the second distribution test.  Its two endpoint bounds are
separate assumptions: throughout this section, \(\sigma\) is a law
on units, \(\nu\) is a law on points, and
\begin{equation}
\label{eq:asym-caps}
 \sigma\le 2^{3000gN}\pi^2,
 \qquad \sigma_1,\sigma_2\le L\pi,
 \qquad \nu\le L\pi.
\end{equation}
Here \(\pi\) is the thinned law and \(\rawlaw\) is the original
frame law.  In particular only their primal marginals agree.  We
use \(\varepsilon_{\rm ch}=10|\cE|2^h\rho\) as an upper bound
on the exponent lost in comparing \(\pi^2\) with \(\rawlaw^2\).
All restrictions made below have fixed positive mass, so all their
constant marginal factors can be specified before the channel size
is chosen.

Draw \(A\sim\sigma\) and, independently, two points
\(B_1,B_2\sim\nu\).  We use this independent experiment for all
calculations involving the two points.  By
\Cref{cor:independent-point-leaf},
\begin{equation}
\label{eq:asym-internal-hole}
 \PP\{B_1B_2\text{ is a hole}\}
 \le O_g\bigl(L^2 2^{2\dim\cB-N}\bigr)=o(1).
\end{equation}
Indeed the two primal spans in each component are disjoint except
on the displayed exceptional mass.  On their complement, the first
hole equation forces both profiles to be zero, contradicting the
last hole equation.  Only after the independent-point calculations
will we delete these internal holes.  The resulting unit law has a
single unpinned leaf and
\begin{equation}
\label{eq:asym-zero-space}
 C_{B,1}=C_{B,2}=0.
\end{equation}
The full-frame image bounds needed for that leaf pay
\(\varepsilon_{\rm ch}N\); they do not assert an unchanged
full-frame marginal.  Conditioning on a nonhole costs \(1+o(1)\).

\subsection{No common prediction for the point side}

Write \(\Omega_{l,q,n}\) for the single-key reference space from
\Cref{sec:preparation}.  For a channel frame \(v\) and a key
\(k=((x_e,y_e))_{e\ni l}\) in this space, put
\[
 \chi_v(k)=(-1)^{u_v(k)},\qquad
 u_v(k)=\sum_{e\ni l}
       (Y_{v,e}\trans x_e)\cdot(X_{v,e}\trans y_e).
\]
Every component has \(x_e\cdot y_e=q\).  This notation agrees
with evaluating \(u_v\) on the tensor represented by the key.

\begin{lemma}[Channel correlations on a common key]
\label{lem:asym-channel-orthogonality}
Let \(v,w\) be two frames for which the concatenated channel
columns \([X_{v,e},X_{w,e}]\) and
\([Y_{v,e},Y_{w,e}]\) have rank \(2h\), in every component
of the star of \(l\).  Uniformly in these frames and in \(q\),
\begin{align}
 \E_{\Omega_{l,q,n}}\chi_v
   &=2^{-h(g-1)}+O_g(2^{-N+4h}),\label{eq:asym-channel-mean}\\
 \E_{\Omega_{l,q,n}}\chi_v\chi_w
   &=2^{-2h(g-1)}+O_g(2^{-N+4h}).
   \label{eq:asym-channel-correlation}
\end{align}
For independent frames of law \(\nu\le L\pi\), these rank
conditions fail with probability \(o(1)\).
\end{lemma}

\begin{proof}
Before imposing \(x\cdot y=q\) in one component, the four
channel evaluation vectors
\[
 (Y_v\trans x,Y_w\trans x,X_v\trans y,X_w\trans y)
\]
are uniform on \(\F^{4h}\).  Conditional on their values,
\(x,y\) are independent uniforms on affine subspaces of
codimension \(2h\).  The pairing between their translation
spaces has rank at least \(N-4h\).  For a bilinear form of
rank \(R\), its sign average on an affine product is either
zero or has absolute value \(2^{-R}\): average first in one
variable, leaving at most a \(2^{-R}\) fraction of the other
variable.  Thus every fixed evaluation tuple has diagonal
probability \(\frac12+O(2^{-N+4h})\).  Conditioning on the
diagonal changes its uniform law by \(O(2^{-N+4h})\) in total
variation.  Deleting zero key vectors incurs an additional
\(O(2^{-N})\) error.

For independent uniform bits \(a,b\),
\(\E(-1)^{ab}=1/2\).  There are \(h\) such products for one
frame and \(2h\) independent products for the pair of frames
in each of the \(g-1\) components.  Independence across
components proves the two formulas.

Under \(\rawlaw^2\), two independent channel frames have
concatenation rank less than \(2h\) with probability at most
\(O_g(2^{2h-N})\), by exposing their uniform injective columns.
The domination
\(\nu^2\le L^2 2^{\varepsilon_{\rm ch}N}\rawlaw^2\)
still makes this probability exponentially small, since
\(\varepsilon_{\rm ch}<1/2\).  This proves the last assertion.
\end{proof}

\begin{lemma}[Absence of a point-side prediction]
\label{lem:asym-no-point-prediction}
In the common limit of the independent two-point law, there is
no positive mass on which, at either endpoint, a prediction
\begin{equation}
\label{eq:asym-point-prediction}
 p_{B,j}=\alpha h_{B,j}+f(l,q,\xi)
\end{equation}
holds for a fixed measurable common-signature function \(f\),
a fixed bit \(\alpha\), and all surviving unary types after
bounded label exclusions.  The assertion holds for every fixed
positive mass, however small, and after deletion of internal holes.
\end{lemma}

\begin{proof}
Suppose the limiting mass is positive.  Approximate \(f\), on
each of the finitely many tag/diagonal reference spaces, by binary
cylinder functions.  For a fixed approximation, the maximum
prediction error over the surviving unary types is a continuous
function of the finitely many stored status-cylinder coordinates.
Minimizing over the finite allowed exclusion choices preserves
continuity.  The strict-error transfer in the proof of
\Cref{lem:status-global-prediction} therefore gives a subsequence
and global key functions \(f_n\) for which the finite prediction
error tends to zero on a set of pairs of mass at least some
\(\delta>0\).  The functions are fixed before sampling these
pairs.  We use no supremum over functions chosen by a sampled
orientation.

There is a fixed surviving generic label on positive pair mass:
average over the finite label set and retain one of its positive
terms.  Fix also a tag, a feasible diagonal, and a feasible generic
tester summary \(\tau\).  This summary has a fixed positive
nominal probability and fixes the role \(a(\tau)\).  Conditioning
the vanishing error on it still leaves vanishing error.  Absorb
the term \(\alpha a(\tau)\) into the global binary test
\(\phi_n(k)\).  Write the two points in the relevant order as
\(U,V\), so that \(p_{B,j}\) evaluates the channel form of
\(V\) on the primal key supplied by \(U\).  Set
\[
 \begin{split}
 D_n(U,V)&=\PP_z\{u_V(K_U(z))\ne\phi_n(K_U(z))
                       \mid \tau\},\\
 d_n(V)&=\PP_{k\sim\Omega_{l,q,n}}
                         \{u_V(k)\ne\phi_n(k)\}.
 \end{split}
\]
For each fixed \(V\), the one-endpoint assertion of
\Cref{lem:typicality} applies with that entire frame as external
data.  Its exceptional set depends on this prescribed test, and
its bound is uniform in the external data.  The raw mean is
\(d_n(V)+o(1)\) by the single-key frame law.  The marginal
of the primal frame of \(U\) is at most \(L\) times its
\(\rawlaw\)-marginal.  Consequently, for every fixed
\(t>0\),
\[
 \PP_{U,V\sim\nu}\{|D_n(U,V)-d_n(V)|>t\}=o(1).
\]
For example the variance proof bounds the left side, apart from
the negligible reference error, by \(LC2^{-cn}/t^2\).
Choose \(t_n\downarrow0\) slowly enough for the exceptional
probability still to vanish.  If \(D_n\le\epsilon_n\to0\)
on the successful pair set, its projection to the \(V\)-coordinate
therefore yields
\[
 G_n=\{v:d_n(v)\le\epsilon_n+t_n\},\qquad
 \nu(G_n)\ge\delta-o(1).
\]
Any adaptive choice of the successful pair set is harmless here:
the typicality exception was bounded on the independent product
law before restricting it.

Take two independent frames \(V,W\) from \(\nu\) conditioned
on \(G_n\).  The probability of channel rank failure is at most
its unconditioned \(o(1)\) probability divided by
\(\nu(G_n)^2\), and hence still tends to zero.  On a pair
with full ranks, agreement of both channel signs with the same
test gives
\[
 \E_{\Omega_{l,q,n}}\chi_V\chi_W
 \ge 1-4(\epsilon_n+t_n).
\]
This contradicts \eqref{eq:asym-channel-correlation}, whose
right side tends to \(2^{-2h(g-1)}<1\).  This use of two
successful frames is essential: it introduces no required lower
bound on \(\delta\) besides positivity.  Finally,
\eqref{eq:asym-internal-hole} changes all bounded error or event
probabilities by \(o(1)\), so deleting internal holes cannot
create the asserted limiting prediction.
\end{proof}

\subsection{A singleton phase estimate}

The point side needs only one phase calculation.  In the following
lemma the marks on \(A\) belong to its first effective spaces,
and \(\Delta_A=U_{A,1}\lambda_{A,1}+U_{A,2}\lambda_{A,2}\).

\begin{lemma}[One point, then two independent points]
\label{lem:asym-singleton-phase}
Let a marked unit law have a fixed marginal bound relative to
\(\pi\) and an absolute joint bound \(2^{(D+o(1))N}\)
relative to \(\rawlaw^2\).  Suppose \(\Delta_A\ne0\)
throughout, with its component rank bounded by \(r\), and fix
\(c\in\F\).  One can retain a fixed positive fraction of
the unit law such that, for an independent point \(Z\sim\nu\),
\[
 \PP\{u_Z(\Delta_A)=c\text{ and }R(A,Z)\}>.49.
\]
Here \(R\) is either the whole space or the transpose-rank
condition on the offset list used below.  The retained fraction
can be fixed before \(h\); the sufficiently large choice of
\(h\) may depend on the fixed marginal constants.  With two
independent points, under this same retained unit law,
\begin{equation}
\label{eq:asym-two-phase-bits}
 \PP\{u_{B_1}(\Delta_A)=u_{B_2}(\Delta_A)=c,
                  \ R(A,B_1),R(A,B_2)\}>.49^2.
\end{equation}
\end{lemma}

\begin{proof}
Use the cover threshold \(p_0\), even moment \(s\), and
rank-growth constants in \Cref{lem:no-cover-rank} and the proof
of \Cref{thm:phase-alternative}.  Explicitly, \(s\) is the
largest power of two at most \(c_sN\), where
\(c_s=10^{-12}(1+r)^{-4}\), so \(s\ge c_sN/2\) for large
\(n\).  Put \(\theta=10^{-4}\) and
\(\Gamma=D+2r+10\).  If no cover of the specified
size has unit mass at least \(p_0\), expand the moment against
independent uniform channel columns of the one point:
\[
 \E_{\rm iid\ channels}
     \left|\E_A(-1)^{u_Z(\Delta_A)}\right|^s
 \le \E_{A_1,\ldots,A_s}
       2^{-h\rank(\Delta_{A_1}+\cdots+\Delta_{A_s})}
 \le 4^s p_0^{q_0s/2}+2^{-hq_0s/4}.
\]
This is \eqref{eq:phase-moment} with no reciprocal phase and no
candidate value of \(\Delta_Z\).  The inequalities
\eqref{eq:early-p0} and \eqref{eq:early-h-moment}, followed
by Markov's inequality at threshold \(\theta\), bound the
raw channel exception by \(C_h2^{-\Gamma N}\).  The channel
comparison in \Cref{lem:frame-images} is included in the fixed
factor \(C_h\).  Paying \(L2^{\varepsilon_{\rm ch}N}\)
once leaves an exponentially small point-law exception.  Thus
\(\left|\E_{A,Z}(-1)^{u_Z(\Delta_A)}\right|<.01\)
for large \(n\).  Take \(R\) to be the whole space; each
value of the bit has probability greater than \(.495\).

Otherwise retain the \(p_0\) mass in a fixed cover and delete
the negligible set whose individual primal spans meet that
cover.  The deletion uses the unchanged primal marginal.  The
lifting identity \eqref{eq:cover-lifting} gives offset records
with at most \(2^{.002N}\) possibilities and offset basis size
\(1\le t\le2r\).  At one fixed record, let \(Z_A\) be its
companion tuple.  Its unnormalized subprobability law satisfies
\[
 p_{\max}(Z_A)\le2^{-tN+.01N}
\]
by \eqref{eq:companion-pmax}.  For that record let \(R(A,Z)\)
require the channel transposes of \(Z\) to have full rank on
the plus and minus offset lists.  The transpose estimate in
\Cref{lem:channel-thinning}, multiplied by the fixed point
marginal factor, gives
\[
 \PP\{R(A,Z)\text{ fails}\}
 \le L\,O(r)2^{2r-h}+o(1)<.01
\]
after increasing \(h\).  We may, and do, make this bound less
than \(.005\).

On the retained rank event the output tuple
\(\Phi_{u_Z}(\omega_A)\) has point masses at most
\(2^{-tN+.01N}\).  To check this directly, condition on the
point's primal frame and freeze the bounded lists of transpose
coefficients on the offsets.  Each list has full rank.  Ignoring
the equations defining those coefficients, their prescribed
images under the opposite channel maps cost
\(t(N-\dim\cB)\) bits.  The number of coefficient lists is
fixed in \(n\).  The density loss from thinning and the ratio
\(\dim\cB/N\) are absorbed into \(.01N\).  This is the
singleton image calculation of \eqref{eq:singleton-joint-pmax}
with no point-side companion tuple.

The phase is \(Z_A\cdot\Phi_{u_Z}(\omega_A)\), plus a
function of the point alone at this fixed record.  The Walsh
bound \Cref{lem:walsh} therefore bounds its unnormalized signed
mass, including the rank filter, by
\[
 2^{tN/2}
 \left(2^{-tN+.01N}2^{-tN+.01N}\right)^{1/2}
 =2^{(-t/2+.01)N}.
\]
Summing over records gives at most \(2^{-.488N}\).  Since
\(\PP(R)> .995-o(1)\), the mass where \(R\) holds and
the phase bit has either prescribed value exceeds \(.49\).
This proves the singleton assertion.

For the last assertion put
\[
 f(A)=\PP_{Z\sim\nu}\{u_Z(\Delta_A)=c, R(A,Z)\mid A\}.
\]
The two points are conditionally independent given \(A\), with
the identical separate rank filter, so the left side of
\eqref{eq:asym-two-phase-bits} is
\[
 \E_A f(A)^2\ \ge\ (\E_A f(A))^2\ >\ .49^2.
\]
This argument asserts an averaged singleton estimate; it does
not require a pointwise estimate for each unit.  Internal holes
have not yet been deleted.
\end{proof}

\subsection{The mixed dual relations and the conflict bound}

\begin{proposition}[The asymmetric distribution test]
\label{prop:asymmetric-distribution}
For every fixed \(L\), with the late parameters chosen as above,
all sufficiently large \(n\) and all laws satisfying
\eqref{eq:asym-caps} obey
\[
 \PP_{A\sim\sigma,\ Z\sim\nu}
       \{A\text{ and }Z\text{ conflict}\}\ge2^{-100gN}.
\]
Thus \Cref{thm:distribution-tests}\textup{(ii)} holds.
\end{proposition}

\begin{proof}
Suppose a violating sequence exists.  Peel the \(A\) law once
by \Cref{lem:peeling}\textup{(a)}; retain the actual leaf weights
and the resulting first effective spaces.  After making all
restrictions on this law, we will apply part \textup{(b)} once
to restore the required normalized-leaf image bounds without
adding pins.  The margins between \(3000g\) and
\(D=4000g\) absorb thinning and fixed normalizations.
Throughout, the mixed marginal bounds used in
\Cref{sec:histograms,sec:common-limit} concern the primal frames.

\paragraph{Preparing the unit and the point pair.}
Make the following subsequential choice on the \(A\) law.
If an empty prediction exists on mass at least \(10^{-6}\),
retain it.  Here empty means \(\alpha_1=\alpha_2=1\) and
the offsets are \(\gamma_1q,\gamma_2q\) with fixed bits
\(\gamma_i\), modulo errors tending to zero in the reference
spaces.  \Cref{lem:empty-status} exactifies it, after negligible
deletion, to
\begin{equation}
\label{eq:asym-exact-A}
 p_{A,i}=T(\lambda_{A,i},\cdot)+a
                  \quad\hbox{on all of }\cX,\qquad i=1,2.
\end{equation}
In particular its proof eliminates the multiples of \(q\),
using the odd number of star atoms.  Fix
\(c=a(\lambda_{A,1})+a(\lambda_{A,2})\) and the indicator
of \(\Delta_A=0\), losing a factor at most four.  If no
such empty prediction exists, retain the first law without
prediction preparation.  Only this specialized empty prediction
is tested; arbitrary common predictions need not be excluded.

In the predicted branch with \(\Delta_A\ne0\), apply
\Cref{lem:asym-singleton-phase}.  In the branch with
\(\Delta_A=0\), necessarily \(c=0\).  Indeed the identities
\eqref{eq:asym-exact-A} give
\((u_{A,1}+u_{A,2})U_{A,i}=T(\lambda_{A,i},\cdot)+a\),
since \(u_{A,i}U_{A,i}=0\); if \(c=1\), these profiles and
\(\Delta_A=0\) satisfy all hole equations for the unit \(A\),
a contradiction.  Thus in either predicted branch the independent
experiment has fixed positive mass on
\begin{equation}
\label{eq:asym-compatible}
 u_{B,j}(\Delta_A)=c\qquad(j=1,2).
\end{equation}
In the nonzero case this mass exceeds \(.49^2\) under the
normalized retained \(A\) law.  Before constructing either
common model, apply \Cref{lem:peeling}\textup{(b)} to the
combined restriction made on \(A\).  More explicitly, let
\(w_\lambda\) be the original first-leaf weights and let
\(q_\lambda\) be the fraction retained in leaf \(\lambda\)
by all the \(A\)-only restrictions just made.  Their total
mass \(q=\sum_\lambda w_\lambda q_\lambda\) is bounded below
by a fixed positive constant.  Discard the leaves with
\(q_\lambda<2^{-\zeta N}\).  They have mass at most
\(2^{-\zeta N}/q=o(1)\) in the restricted law; every remaining
normalized leaf has the image bound
\eqref{eq:final-leaf-minentropy} and an absolute density
\(2^{O(N)}\) relative to \(\rawlaw^2\).  The pins and first
effective spaces have not changed, and the actual remaining
leaf weights are retained.  This operation depends only on
\(A\), so the two points remain independent conditional on
\(A\).  On the predicted branch the two-bit compatibility mass
is still at least \(.49^2-o(1)\).

Now delete internal \(B\) holes, costing \(o(1)\) by
\eqref{eq:asym-internal-hole}, and take the single unpinned
\(B\) leaf from \Cref{cor:independent-point-leaf}.
Apply the independent-law version of \Cref{lem:injection} to
these two laws, with failure probability less than \(.01\).
On the predicted branch compatibility and injection therefore
coexist on mass at least \(.49^2-.01-o(1)>.22\), for large
\(n\).  On the branch without prediction a passed injecting
actual table has mass at least \(.99\).  These are subtractions
of exceptional mass, and require no independence between the
exceptional events and compatibility.

\paragraph{Separate flags for candidate tables.}
We record this mass in the permitted finite-table format.
For each leaf pair enumerate all numerical candidate small
tables \(T\), as in \Cref{def:limit-option}; there are finitely
many, with a bound independent of \(n\).  A fixed \(T\) has
separate admissibility flags \(F_{A,T}(A)\) and \(F_{B,T}(B)\)
by \Cref{lem:small-tables}.  In the predicted branch define
\begin{equation}
\label{eq:asym-table-flag}
 H_{A,T}(A)=\1\left\{
   \sum_{i=1}^2(u_{B,j}U_{A,i})^*_T(\lambda_{A,i})=c
       \quad(j=1,2)\right\}.
\end{equation}
The star denotes contraction in the numerical table \(T\),
as in \eqref{eq:status-table-delta}.  It is calculated from
the entries of \(T\), the recorded tensor coordinates of the
\(C_{A,i}\) bases, and the bounded coordinate vectors of the
marks \(\lambda_{A,i}\).  Thus, for a fixed leaf pair and
\(T\), \(H_{A,T}\) is a flag of \(A\) alone.  Include the
mark coordinates among the finite orientation metadata and
replace \(F_{A,T}\) by \(F_{A,T}H_{A,T}\).  In the branch
without prediction set \(H_{A,T}=1\).  Injection is a finite
matrix-rank test on the candidate table and the leaf metadata.

Each pair counted by the ambient compatibility and injection
estimate has at least one passed candidate: take \(T\) to be
its actual small table.  Therefore the union of these permitted
candidate-table events has at least the mass just established.
We store their entire finite flag vector in the common model.
This argument does not impose equality of a candidate table
with the actual cross entries between private directions;
such an equality would be a joint condition.  The actual table
is used only to witness the existence of a passed candidate.
Nor do we condition the two unit laws on the ambient phase
event \eqref{eq:asym-compatible}.  They remain independent
laws with the separate flags just described.

\paragraph{Solving the mixed scalar system.}
Construct the common model for these two different unit laws,
retaining the ordered \(A\)--\(B\) type filter.  The dual
relations \eqref{eq:status-dual-relation} become
\begin{equation}
\label{eq:asym-dual}
 s_j(p_{A,i}+h_{A,i})+r_i(p_{B,j}+h_{B,j})
 =\mathbf c_{ij}\cdot\mathbf x_{A,i}
       +\beta_{ij}(h_{A,i}+h_{B,j})+\gamma_{ij}q,
\end{equation}
because \(C_{B,j}=0\).  We first show that no relation with
some \(r_i=1\) can occur on positive model mass, for any of
the permitted bounded exclusions.  If it did, Fubini fixes an
\(A\) array with a positive-mass set of such \(B\) arrays.
Pigeonhole its finite coefficient and exclusion data and fix a
surviving generic label on the \(A\) side.  At a fixed common
signature, the two status distributions are independent.  The
two expressions in \eqref{eq:asym-dual} must agree almost
surely, because a mismatching positive-density pair would be an
available generator violating the relation.  Two independent
bits which agree almost surely are both deterministic.  The
generic \(A\) type has total status density one on every
signature space, so its deterministic value gives a fixed
common function on all surviving \(B\) types.  Solving
\eqref{eq:asym-dual} for \(p_{B,j}\) gives
\[
 p_{B,j}=(1+\beta_{ij})h_{B,j}+f(l,q,\xi).
\]
The strict-error cylinder transfer used in
\Cref{lem:asym-no-point-prediction} applies, and that lemma
rules this out.  The finite number of possible dual and
exclusion records implies that, almost surely, every remaining
dual relation has \(r_1=r_2=0\).

In the branch without an empty prediction, suppose obstructable
pairs had probability at least \(.1\).  A nonzero weight now
has some \(s_j=1\).  Generic positivity of the two possible
\(B\)-role values in \eqref{eq:asym-dual} forces
\(\beta_{ij}=0\).  The same relation, for both \(i\), is
then
\begin{equation}
\label{eq:asym-empty-extraction}
 p_{A,i}=T(\lambda_{A,i},\mathrm{atom})
                         +h_{A,i}+\gamma_{ij}q.
\end{equation}
It holds on all surviving unary types, by comparison to a fixed
generic \(B\) type of total density one.  These are empty
predictions.  There are only two choices of \(j\) and four
offset pairs \((\gamma_{1j}q,\gamma_{2j}q)\); therefore
one choice occurs on \(A\) mass at least \(.1/8\).
The marked profiles and exclusions may still vary with \(A\),
as allowed in \Cref{def:prediction}; their number is not a
further pigeonhole cost.  Cylinder transfer with strict error
events, as in \Cref{lem:status-global-prediction}, loses at
most a factor two in this mass and supplies a finite empty
prediction of mass much larger than \(10^{-6}\).  This is a
contradiction.  Thus obstructable mass is less than \(.1\),
whereas passed injecting tables have mass at least \(.99\).
Their intersection with nonobstructable pairs has mass greater
than \(.89\).  On it \Cref{lem:status-duality} solves the
scalar system for every such table and every allowed exclusion.

On the predicted branch, substitute \eqref{eq:asym-exact-A}
into \eqref{eq:asym-dual}.  The full generic base projection
forces
\[
 \mathbf c_{ij}=s_j\boldsymbol\lambda_{A,i},
 \qquad \beta_{ij}=0,\qquad\gamma_{ij}=0.
\]
For \(s_j=0\) all these coefficients are zero.  Fix any
passed injecting candidate table \(T\), including the flag
\(H_{A,T}=1\).  At its local targets the dual evaluation is
\[
 \sum_{i,j}\mathbf c_{ij}\cdot\mathbf t_{A,i;j}
 =\sum_j s_j\left(
       \sum_i a(\lambda_{A,i})
       +\sum_i(u_{B,j}U_{A,i})^*_T(\lambda_{A,i})\right)=0
\]
by \eqref{eq:asym-table-flag} and
\(\sum_i a(\lambda_{A,i})=c\).  This calculation uses the
selected candidate's starred contractions throughout; it does
not require that candidate to equal the actual cross table of
the sampled pair.  Every dual relation vanishes on its target.
By \Cref{lem:status-duality}, the scalar system is soluble on
the positive mass having a passed compatible injecting
candidate.  The reasoning uses only surviving generic labels
and hence persists through the further bounded exclusions for
fresh lists.

\paragraph{Positive overlap and collision.}
In either branch we now have the hypothesis of the mixed,
filtered \Cref{prop:positive-overlap}.  Equivalently,
\Cref{lem:fresh-lists} supplies at most seven generators per
cross edge, with distinct labels at each endpoint and positive
status overlap at every matched position.  The scalar recipe
and gradient realization produce a usable numerical option;
the filtered zero-product conclusion of
\Cref{prop:limit-overlap} cannot hold on this positive mass.
Consequently some option has a positive subsequential expected
unnormalized overlap.  This positive constant need not be
uniform over all possible limit models.

The collision criterion \Cref{prop:collision} therefore gives
four cross holes between \(A\) and \(B\) with probability
\[
 2^{-(k_{\max}+.03+o(1))N}.
\]
All retained-law normalizations have fixed positive mass and
are included in the \(o(1)N\) loss.  In particular these four
holes imply that \(A\) conflicts with \(B_1\).  Viewed in
the original independent experiment, the latter has exactly
the probability of a conflict with one point of law \(\nu\).
Since
\[
 k_{\max}+.03=56(g-1)+.03<100g,
\]
this contradicts the assumed bound below \(2^{-100gN}\) for
all sufficiently large indices of the violating sequence.
The asserted uniform eventual distribution test follows.
\end{proof}
